% Microarchitecture chapter for Simple 8-Bit Microprocessor
\section{Microarchitecture}
This chapter describes the microarchitecture of the Simple 8-bit CPU and provides a visual datapath diagram.

\subsection{High-level description}
The CPU is a single-cycle, accumulator-based design. There is no pipelining: on each rising clock edge, the PC, ACC and HALT state are updated. Instruction fetch reads IMEM[PC] into IR; the opcode is decoded in combinational logic to drive the datapath and control signals. Memory writes to DMEM occur in the sequential stage when a store is executed.

\subsection{Datapath and control}
The datapath components are:
\begin{itemize}
  \item Program Counter (PC) -- 8-bit register (low 4 bits used for addressing internal memory).
  \item Instruction Memory (IMEM) -- 16 x 8-bit ROM/array.
  \item Instruction Register (IR) -- holds fetched instruction.
  \item Decoder / Control Logic -- combinational logic deriving control signals from IR[7:4].
  \item Accumulator (ACC) -- 8-bit register used by ALU operations.
  \item ALU -- supports ADD and SUB; combinational, 8-bit result feeds next ACC.
  \item Data Memory (DMEM) -- 16 x 8-bit array with write enable.
\end{itemize}

\subsection{TikZ Datapath}
The following diagram shows the datapath and control flow.

\begin{figure}[!ht]
\centering
\begin{tikzpicture}[>=latex, node distance=8mm, thick, every node/.style={font=\sffamily}]
% Boxes
\node[draw, rectangle, minimum width=22mm, minimum height=8mm] (PC) {PC [7:0]};
\node[draw, rectangle, right=20mm of PC, minimum width=28mm, minimum height=8mm] (IMEM) {IMEM};
\node[draw, rectangle, below=8mm of IMEM, minimum width=28mm, minimum height=8mm] (IR) {IR};
\node[draw, rectangle, right=20mm of IR, minimum width=28mm, minimum height=8mm] (DEC) {Decoder / Ctrl};
\node[draw, rectangle, above=16mm of DEC, minimum width=28mm, minimum height=8mm] (ALU) {ALU};
\node[draw, rectangle, above=16mm of ALU, minimum width=28mm, minimum height=8mm] (ACC) {ACC [7:0]};
\node[draw, rectangle, right=20mm of DEC, minimum width=28mm, minimum height=8mm] (DMEM) {DMEM};

% Connections
\draw[->] (PC) -- node[midway, above] {addr[3:0]} (IMEM);
\draw[->] (IMEM) -- node[midway, right] {instr} (IR);
\draw[->] (IR) -- node[midway, above] {opcode/operand} (DEC);
\draw[->] (DEC) -- ++(6mm,0) |- (ALU) node[pos=0.25,right]{op};
\draw[->] (ACC) -- node[midway, left]{acc} (ALU);
\draw[->] (ALU) -- node[midway, right]{result} (ACC);
\draw[->] (DEC) -- ++(6mm,0) |- (DMEM) node[pos=0.2,right]{mem addr/wr};
\draw[->] (ACC) -| node[pos=0.25,below]{data} (DMEM);
\draw[->] (DMEM) -| ++(-18mm,-14mm) |- (ACC) node[pos=0.9,left]{ld};

% Control signals
\node[below=6mm of DEC] (ctrlnote) {Control signals: mem\_wr\_en, next\_pc, next\_acc, halt};

\end{tikzpicture}
\caption{Microarchitecture datapath (simplified).}
\label{fig:microarch}
\end{figure}

\subsection{Frequency and area estimation}
A generic Yosys synthesis is available for fast structural inspection, while the
repository also contains a complete SKY130 implementation flow. Results from
one checked-in flow run are summarized below.

\subsection{Synthesis results (summary)}
\begin{itemize}
  \item Generic Yosys hierarchy: \textbf{290 cells}, including generic logic,
        registers, and memory primitives.
  \item SKY130HD final report: \textbf{15,142\,\textmu m\textsuperscript{2}}
        placed cell area at a 10\,ns clock constraint.
  \item SKY130HD core die: approximately \textbf{193.255\,\textmu m} square
        at the selected 35\% initial utilization.
  \item The final routed layout contains 4,820 instances including standard
        cells, taps, clock/timing buffers, and filler cells.
\end{itemize}

\subsection{Assumptions and notes}
The physical numbers are flow results, not universal limits for the CPU. They
depend on the SKY130HD library, constraints, placement density, tool versions,
and the exposed memory-loader ports. The 10\,ns clock is a design target; it
must not be interpreted as a silicon frequency guarantee. The final GDSII and
complete reports are generated by \texttt{make gds}; the stable deliverable is
\texttt{outputs/cpu\_sky130hd.gds}.

See \texttt{artifacts/} for the generic Yosys netlist and
\texttt{build/orfs/} for the OpenROAD reports and intermediate databases.

\subsection{Discussion}
The critical logic is expected through the instruction decoder and ALU feeding
the accumulator. The small memories are inferred from RTL and mapped to
standard-cell logic by the open-source flow; a production design could replace
them with dedicated SKY130 memory macros and add a pad/power wrapper.

% End of microarchitecture chapter
